\section{Supporting Analyses and Stopped Routes}
\label{app:secondary}

This appendix separates secondary evidence from work that did not reach an interpretable scientific test.
A working pipeline, a restricted local diagnostic, and a stopped route answer different questions.
Appendix~\ref{app:evidence-provenance} states how far each record's conclusion may be taken, and Appendix~\ref{app:sensitivities} reports the full numerical results and sensitivity analyses.
The headline empirical verdicts remain in Section~\ref{sec:results}.

\subsection{Pipeline validation and scientific scope}

The end-to-end pipeline was first tested in a setting where the expected signal was known.
Real Pereira sentences \autocite{pereira2018} were paired with synthetic responses generated from teacher features, named nuisance variables, and noise.
The test exercised ordinary and target-guided \ac{kd}, the temporary target head, middle-layer extraction, the contiguous held-out split, fold-local ridge fitting, nuisance subtraction, and aggregation across optimization seeds.
The expected teacher--student ordering appeared, the nuisance component remained small, and excessive target weight degraded the result rather than receiving an automatic advantage. \evd{E001}

These observations rule out several basic implementation failures.
They show that the held-out evaluation tail was not used for target-guided training, that the scoring pipeline could recover a planted feature--response relation, and that the analysis did not reward every target-guided run by construction.
They do not establish that an \ac{lm} predicts recorded brain responses, that the nuisance set is sufficient for real stimuli, or that a brain-response objective improves a student.
The synthetic test therefore validates the end-to-end implementation, not the scientific claim.

\subsection{Response averaging and recorded-response robustness}

Two averaging analyses addressed different questions.
The first changed how many participants contributed to the recorded-response intervention target.
Because each target was rebuilt from a different participant subset, increasing the subset size also changed participant identity, response reliability, and target construction.
A corrected analysis showed that the apparent gain in controlled brain predictivity did not rise steadily with subset size.
The result therefore does not identify a causal benefit of adding more participants to the training target. \evd{E010}

\looseness=1 The second analysis held the frozen representation and encoding procedure fixed while averaging increasing numbers of recorded responses.
Controlled brain predictivity increased as averaging suppressed idiosyncratic and measurement noise and emphasized the shared stimulus-evoked component.
This is a valid signal-to-noise gain for a group-response estimand.
It changes what is being predicted, however, and does not supply participant replication or rescue the individual-response~\mbox{intervention}. \evd{E014}

The remaining analyses tested whether specific design choices concealed a recorded-response benefit.
Table~\ref{tab:recorded-robustness} maps each change to its objection and scope.

\begin{table}[tbp]
  \centering
  \footnotesize
  \setlength{\tabcolsep}{4pt}
  \renewcommand{\arraystretch}{1.04}
  \caption{Recorded-response robustness analyses and their scope.}
  \label{tab:recorded-robustness}
  \begin{tabularx}{\textwidth}{
    @{}
    >{\raggedright\arraybackslash}p{29mm}
    >{\raggedright\arraybackslash}p{38mm}
    >{\raggedright\arraybackslash}X
    >{\raggedright\arraybackslash}p{39mm}
    @{}
  }
    \toprule
    Objection & Design change & Supported disposition & Remaining open \\
    \midrule

    Auxiliary weight was poorly chosen.
    & Repeated the recorded-response intervention over the tested weight range.
    & Only the highest tested weight gave a fold-level interval against the matched permuted control that excluded zero, and at a substantial language-quality cost, so no quality-matched target-specific gain is supported.\evd{E005}
    & Other objectives or schedules could behave differently. \\

    \specialrule{0.15pt}{0.3ex}{0.3ex}

    Adapter capacity was too limited.
    & Increased the tested low-rank adaptation capacity.
    & The tested capacity increase did not recover the effect.\evd{E011}
    & This does not test every stronger manipulation. \\

    \specialrule{0.15pt}{0.3ex}{0.3ex}

    Pointwise regression was the wrong loss.
    & Replaced it with a contrastive response objective.
    & The alternative loss did not produce a reliable target-specific benefit.\evd{E013}
    & Other response objectives remain possible. \\

    \specialrule{0.15pt}{0.3ex}{0.3ex}

    Regional targets hid a voxelwise effect.
    & Tested the available naturalistic voxelwise training procedure.
    & Target-matched tuning was neutral at the lowest tested target weight and degraded held-out predictivity at higher weights, never exceeding either the untuned model or its permuted twin, so the intended change was never induced.\evd{E013}
    & The broader multi-participant full-parameter route was not built. \\

    \specialrule{0.15pt}{0.3ex}{0.3ex}

    Parameter-efficient tuning was too restrictive.
    & Tested full-parameter tuning, including a gentler three-participant probe.
    & Aggressive tuning caused forgetting; the gentler probe showed no reliable target-specific gain.\evd{E017}
    & Three participants do not support population inference. \\

    \bottomrule
  \end{tabularx}
\end{table}

These analyses weaken the listed specific alternative explanations.
They do not establish that every recorded-response objective must fail.
Section~\ref{sec:results-recorded} states the participant-level intervention conclusion, and Appendix~\ref{app:sensitivities} reports the numerical variants.

\subsection{Reduced external and cross-modal probes}

External protocols were used as local diagnostics only when their decisive ingredients could be approximated with the available data and models.
The standard is therefore narrower than reproduction.
The question is whether a reported effect appears under the local conditions.

The first diagnostic adapted two elements of Negi et al., a differentiable head that aligns model features to response timing and their controls for target specificity \autocite{negi2025}, to the available LeBel responses \autocite{lebel2023} and decoder-only \acp{lm}.
Under the gentlest tested training schedule the student representations that remain after the temporary head is discarded changed little, and stronger schedules degraded language quality severely.
The local runs therefore did not establish a target-specific manipulation over their controls.
The comparison nevertheless differs from the source study in population, response substrate, model family, language, training scale, and control construction.
It is therefore a restricted local failure to establish the target-specific manipulation, not a reproduction or refutation of the published result. \evd{E019}

The second diagnostic adapted a speech-based protocol from Moussa et al. \autocite{moussa2025}.
Starting from the same Wav2Vec2-base model family, the design required a positive phoneme-prediction gain before constructing a matched-quality specificity test.
In the reduced single-participant setup, target loss decreased but the phoneme-gain criterion was not met, so the downstream specificity comparison was not built.
The source study used a larger, multi-participant setting with different timing and model choices.
The result is thus a reduced local non-reproduction of the prerequisite effect, not evidence against the published claim. \evd{E022}

\looseness=1 Neither diagnostic supplied an external positive control for the present intervention pipeline.
They narrow what the local pipeline recovered while leaving the external literature open under its original~\mbox{conditions}.

\subsection{Failed instruments and precondition-stopped routes}

Three routes stopped because an upstream measurement or attribution condition failed.
In each case, continuing to intervention would have produced a number without a defensible interpretation.

The shared-response route asked whether the predictable group response could serve as a ceiling for participant-specific intervention effects.
In principle, this reference estimates the expected response shared across participants for the same stimulus, $m(S)=\mathbb{E}[Y\mid S]$.
In practice, the leave-one-participant reference was not reliable enough, while temporally adjacent samples preserved apparent residual predictivity through leakage and autocorrelation.
Once a clean temporal gap was imposed, the controlled residual predictivity was small but too imprecise to fix a ceiling for participant-specific effects.
The ceiling instrument therefore failed; the analysis does not show that participant-specific structure is absent or that a universal biological ceiling has been reached. \evd{E020}

The reading-time route asked whether a cognitive target could provide a simpler intervention signal than \ac{fmri}.
Training a model with an auxiliary head on residualized reading times initially improved how efficiently a probe on the frozen representation learned held-out text, by more than any control whose target was matched for learnability.
A phase-randomized control preserved the target's marginal distribution and temporal structure, however, and matched that improvement.
The design therefore did not isolate cognition-specific content from target shape.
This attribution failure does not imply that reading time contains no useful signal; it means that the available contrast could not identify that signal. \evd{E021}

The gaze route required three conditions before privileged-information training.
It needed a measurable target, evidence that it added information beyond text, and enough independent groups for leakage-resistant evaluation.
Natural-reading gaze was measurable but did not add information beyond the text itself under the tested assay.
Task-directed gaze was more predictable, but the signal could reflect task intent or label leakage, and the number of independent paragraph groups was too small for a reliable group split.
The intervention therefore stopped at its preconditions.
The route's other planned signal, \ac{eeg}, never reached testing and supplies no evidence. \evd{E024}

Stopping at these failed prerequisites means no downstream number exists that could be misread as an identified intervention effect.
They identify why the proposed instruments could not answer their questions, but they do not turn measurement failure into a biological null.

